\section{Scalar aggregation and matrix query languages}\label{sec:matrix}
We extend the fold argument from binary relations to semiring-valued
matrix computations. We first describe the short three-attribute
expression, then the matrix circuit model, and finally the fold
sensitivity bounds that force exponential arity-elimination overhead.

\Needspace{10\baselineskip}
\subsection{The three-attribute query and its matrix interpretation}\label{sec:weightedG}
Let $\mathfrak S\coloneqq(S,\oplus,\otimes,0_S,1_S)$ be a commutative
semiring and $V$ a finite nonempty common index set, with inputs
$P_i,Q_i:V^2\to S$. Thus $(S,\oplus,0_S)$ and $(S,\otimes,1_S)$
are commutative monoids, $\otimes$ distributes over $\oplus$, and
$0_S$ is absorbing for $\otimes$.
An \emph{annotated relation} assigns an $S$-value to each tuple.
In positive annotated relational algebra, union $\cup$ adds annotations,
natural join $\bowtie$ multiplies them, and projection $\pi_{xy}$ sums
over eliminated attributes while retaining $x,y$; renaming changes
attribute names.
ARA$(3)$ restricts intermediate relations to at most three attributes,
analogous to the ternary assignment relations of \FOthree\
(\S\ref{sec:prelim-logic}). Over Boolean annotations, these operations
are disjunction, conjunction, and existential quantification.

For an explicit ARA$(3)$ expression, let $P_i^{xy}$ denote $P_i$
with attributes $(x,y)$, and similarly for $P_i^{xz}$, $Q_i^{xy}$,
and $Q_i^{xz}$. Our semiring-valued query is
\[
 \mathcal G_k\coloneqq\pi_{xy}\!\left(
 \mathop{\bowtie}_{i=1}^{k}
 \bigl[(P_i^{xy}\bowtie Q_i^{xz})
       \cup(Q_i^{xy}\bowtie P_i^{xz})\bigr]\right).
\]
Every intermediate relation has at most the three attributes
$(x,y,z)$: the $k$ conditions are combined with the \emph{same} $z$, which is
eliminated only at the final projection. Thus the ARA$(3)$ expression
has size $O(k)$. Equivalently, its scalar value is
\[
 \mathcal G_k(x,y)=\bigoplus_{z\in V}\ \bigotimes_{i=1}^{k}
 \bigl[(P_i(x,y)\otimes Q_i(x,z))
        \oplus(Q_i(x,y)\otimes P_i(x,z))\bigr].
\]
For $S=\mathbb B$ with $\oplus=\lor$ and $\otimes=\land$, this
is precisely $G_k(x,y)$, without negation or universal quantification.
For integer min--plus, specializing $\mathfrak S$ to
$(\mathbb Z\cup\{+\infty\},\min,+,+\infty,0)$ gives
\[
 \mathcal G_k(x,y)=\min_{z\in V}\sum_{i=1}^{k}
 \min\{P_i(x,y)+Q_i(x,z),\,Q_i(x,y)+P_i(x,z)\}.
\]

\paragraph{An explicit MATLANG expression.}
Although MATLANG retains only matrices, vectors, and scalars,
the DNF of \S\ref{sec:main-results} gives an expression for
$\mathcal G_k$ using matrix--vector products. For $A\subseteq[k]$,
let $H_A$ be the pointwise semiring analogue of the relation in
\S\ref{sec:main-results}:
\[
 H_A(x,y)\coloneqq
 \bigotimes_{i\in A}P_i(x,y)\ \otimes
 \bigotimes_{i\in A^c}Q_i(x,y),
\]
where an empty product is $1_S$.
Let $\mathbf 1_S$ be the constant column vector with every entry
$1_S$. Computing the column vector $r_A$ and then broadcasting it
along the rows gives
\begin{equation}\label{eq:matlang-dnf}
 r_A\coloneqq H_{A^c}\mathbf 1_S,\qquad
 \mathcal G_k(x,y)=\bigoplus_{A\subseteq[k]}
       H_A(x,y)\otimes r_A(x).
\end{equation}
Distributivity expands the $k$ clauses and interchanges the finite
aggregations over $A$ and $z$, proving the identity for arbitrary
$S$-valued inputs. Each $r_A$ uses one matrix--vector product;
all other steps are entrywise operations or row broadcasting.
Thus $2^k$ such products suffice, even without matrix--matrix
products. Our lower bound shows that an exponential number of
data-index aggregations remains necessary even if both kinds of
products, other permitted operations, and sharing are allowed.

\subsection{Typed matrix circuits}\label{sec:typed}
A semiring is \emph{nontrivial} if $0_S\ne1_S$ and
\emph{idempotent} if $a\oplus a=a$ for every $a\in S$.
In the idempotent case, write semiring addition as $\vee$. The circuit
model below permits arbitrary binary kernels, rather than only
semiring multiplication; it therefore includes the corresponding
MATLANG operations. This strengthening is useful because the lower
bound charges only data-index aggregations.

We use the dimension symbols $n$ and $1$. Intermediate arrays have
shapes $n\times n$, $n\times1$, $1\times n$, or $1\times1$---matrices,
column vectors, row vectors, and scalars. A data-index contraction
of an $a\times n$ and an $n\times c$ array is
\[
 Z(i,j)\coloneqq\bigvee_{z=1}^n\kappa(X(i,z),Y(z,j)),
 \qquad a,c\in\{1,n\},
\]
where the kernel $\kappa:S\times S\to S$ may vary by gate and
need not be associative, monotone, or zero-preserving.
Contractions over an index of size one apply the kernel without
aggregating over data, and are not charged.
We also allow arbitrary uniform entrywise functions: for arrays
$X_1,\ldots,X_r$ of the same shape and a fixed map $f:S^r\to S$,
\[
 F(X_1,\ldots,X_r)(i,j)\coloneqq
 f(X_1(i,j),\ldots,X_r(i,j)).
\]
Transpose exchanges the two indices. Broadcasting a column $\mathbf c$,
a row $\mathbf r$, or a scalar $s$ gives entries $\mathbf c(u)$,
$\mathbf r(v)$, or $s$, respectively, at position $(u,v)$; a scalar
may also be broadcast to a vector. Diagonal extraction returns the
vector with entry $X(u,u)$, and diagonal formation places the vector
entries on the diagonal and $0_S$ off it. We permit uniform constant
vectors and scalars, and a diagonal mask with any two fixed values,
one on and one off the diagonal. These operations are uncharged.
Standard MATLANG operations are described in
\citep{brijder2019tods,brijder2022}.

All circuits have fixed directed acyclic wiring. Functions, kernels,
and uniform scalar constants are fixed before the input is seen,
though they may depend on $n$ and $k$ in a size-specific circuit.
``Uniform'' means that the same entrywise function is applied at
every position of its type and cannot access vertex names.
A kernel receives only its two scalar arguments, not the output
position $(i,j)$, the contracted index $z$, a third matrix entry,
or an encoding of the entire input unless computed by permitted
gates. Data-dependent values and tests are allowed, but there is
no primitive for data-dependent rewiring or arbitrary indexed
access. The input arrays are square $n\times n$ matrices; no order,
vertex identifiers, arbitrary input vectors, or vertex-dependent
constant arrays are supplied for free.

Write $\mvcount C$ for the number of matrix--vector and
vector--matrix data-index contractions and $\mmcount C$ for the
number of matrix--matrix data-index contractions, counting each
distinct gate once. We call these charged gates \emph{aggregation gates}.
All other operations, including vector inner products, are uncharged
in this resource measure. We next show that these
uncharged operations notice no folds, even though some of them may
require compositions when expressed in CoR.

\subsection{Fold sensitivity and finite join breadth}\label{sec:breadth}
Let $(S,\vee,0_S)$ be a join-semilattice with least element $0_S$.
For scalar values $a,b\in S$, use the semilattice order
$a\le b$ if and only if $a\vee b=b$. For min--plus this is the reverse
of the usual numerical order. The semilattice's \emph{join breadth} $\beta$ is the least nonnegative integer such
that every finite join equals the join of at most $\beta$ of its
summands (the empty join is $0_S$). We assume $\beta<\infty$.
For a nontrivial totally ordered semiring with min or max as addition,
$\beta=1$: one summand attains the aggregate. In a product with
$b$ independent coordinates, different summands can determine
the result in different coordinates, and breadth can be $b$.

The significance for folds is immediate. Each fold removes one
color from the hub. If a join is already determined by at most
$\beta$ summands, removing a color outside those summands
cannot change its value. We now make this observation precise
for the three kinds of output pairs involving the hub.

Let $V$ be the common finite index set of the input matrices and
$n\coloneqq|V|$. For color-determined square arrays $X,Y$ and
any kernel $\kappa:S\times S\to S$, put
\[
 (X\agg\kappa Y)(u,v)\coloneqq
 \bigvee_{z\in V}\kappa(X(u,z),Y(z,v)).
\]
Kernels may differ between gates; no order or monotonicity assumption
is made on $\kappa$. Reference arrays are color-determined as in
\S\ref{sec:types}, now with entries $X[\diagc],X[\mu]\in S$.
Their predictions read the same table at the folded coloring
$\chi_\lambda$. Entrywise functions and transpose notice no fold.
For $J\subseteq\Lambda$, define
\begin{align*}
 \kappa_J^{\mathrm L}&\coloneqq\bigvee_{\mu\in J,\,\nu\in\Lambda}
                          \kappa(X[\mu],Y[\nu]),&
 \kappa_J^{\mathrm R}&\coloneqq\bigvee_{\mu\in J,\,\nu\in\Lambda}
                          \kappa(X[\nu],Y[\mu]),\\
 \kappa_J^{\mathrm D}&\coloneqq\bigvee_{\mu\in J}\kappa(X[\mu],Y[\mu]),&
 \kappa_*&\coloneqq\kappa_\Lambda^{\mathrm L}
                     =\kappa_\Lambda^{\mathrm R}.
\end{align*}
The reference product $Z\coloneqq X\agg\kappa Y$ has the table
\begin{align*}
 Z[\diagc]&=\kappa(X[\diagc],Y[\diagc])
            \vee\kappa_\Lambda^{\mathrm D},\\
 Z[\mu]&=\kappa(X[\diagc],Y[\mu])
          \vee\kappa(X[\mu],Y[\diagc])\vee\kappa_*.
\end{align*}
The calculation of \S\ref{sec:responses} remains valid, with
endpoint contributions formed by $\kappa$ and, away from the hub,
outside-witness contributions $\kappa_*$ off the diagonal and
$\kappa_\Lambda^{\mathrm D}$ on it. For the fold of $\lambda$,
put $J\coloneqq\Lambda\setminus\{\lambda\}$. Outside-witness contributions
in the hub row, hub column, and hub diagonal become
$\kappa_J^{\mathrm L}$, $\kappa_J^{\mathrm R}$, and
$\kappa_J^{\mathrm D}$, respectively; their predictions use
$J=\Lambda$. This follows from the same pairs of witness colors,
while idempotence makes their multiplicities irrelevant.

\begin{lemma}[Local bounds for contractions]\label{lem:scalar}
For color-determined square arrays $X,Y$ and an arbitrary kernel
$\kappa:S\times S\to S$, the contraction $X\agg\kappa Y$ notices
at most $3\beta$ folds. If one operand is constant across off-diagonal
colors, it notices at most $\beta$; if both are, it notices none.
For Boolean entries, the general bound improves to two.
\end{lemma}
\begin{proof}
By the calculation above, the gate notices the fold of $\lambda$
only if one of the joins $\kappa_J^{\mathrm L}$,
$\kappa_J^{\mathrm R}$, and $\kappa_J^{\mathrm D}$ with
$J=\Lambda\setminus\{\lambda\}$ differs from its value at
$J=\Lambda$. Choose at most $\beta$ summands of
$\kappa_\Lambda^{\mathrm L}$ whose join is
$\kappa_\Lambda^{\mathrm L}$, and let $K^{\mathrm L}$ be the
set of their row colors, the row color of
$\kappa(X[\mu],Y[\nu])$ being $\mu$. If
$\lambda\notin K^{\mathrm L}$, all chosen summands also occur in
$\kappa_J^{\mathrm L}$, so
$\kappa_\Lambda^{\mathrm L}\le
\kappa_J^{\mathrm L}\le\kappa_\Lambda^{\mathrm L}$.
Define $K^{\mathrm R}$ in the same way from the column colors of
summands of $\kappa_\Lambda^{\mathrm R}$, and $K^{\mathrm D}$
from the colors of summands of $\kappa_\Lambda^{\mathrm D}$.
Hence the gate notices the fold of $\lambda$ only if
$\lambda\in K^{\mathrm L}\cup K^{\mathrm R}\cup K^{\mathrm D}$,
a set of at most $3\beta$ colors.

If $Y[\mu]$ is independent of $\mu$, then
$\kappa_J^{\mathrm L}=\kappa_J^{\mathrm D}$ and
$\kappa_J^{\mathrm R}=\kappa_*$ for every nonempty $J$;
only $K^{\mathrm L}$ is needed. The case of $X$ is symmetric.
If both are constant, all three joins are unchanged.
For Boolean entries, if $\kappa_\Lambda^{\mathrm D}=1$, a single
color $\mu$ with $\kappa(X[\mu],Y[\mu])=1$ supplies a summand
of value one to all three joins whenever $\mu\in J$, so only the
fold of $\mu$ can be noticed. Otherwise the diagonal join is
always zero, and one color for each of the other two joins suffices.
\end{proof}
The displayed reference table is color-determined, so these
calculations also establish the closure required by
Theorem~\ref{thm:test-cost}.

On the reference structure, every vector is constant over its data
index, and every scalar is a single value. Encode a column vector
$\mathbf c$ as the square array $(u,v)\mapsto\mathbf c(u)$,
a row vector $\mathbf r$ as $(u,v)\mapsto\mathbf r(v)$, and a
scalar as a constant square array. Encoded reference vectors are
constant across all edge colors and on the diagonal and are
unchanged by prediction. Diagonal formation and extraction,
broadcasting, and permitted constants notice no folds.
Matrix--vector and vector--matrix contractions notice at most
$\beta$ folds by Lemma~\ref{lem:scalar}; a vector inner product
notices none, since both operands are constant.
Products over an index of size one also notice none.

\begin{lemma}[Cost of typed products]\label{lem:typed-cost}
Every typed matrix circuit $C$ in this model satisfies
\[
 \tcost(C)\le\beta(\mvcount C+3\mmcount C).
\]
For Boolean entries and arbitrary kernels, the bound improves to
$\mvcount C+2\mmcount C$.
\end{lemma}
\begin{proof}
Sum Lemma~\ref{lem:scalar} over the distinct aggregation gates.
All other permitted gates notice none, as shown above.
\end{proof}
\begin{center}
\renewcommand{\arraystretch}{1.1}
\begin{tabular}{p{0.62\linewidth}l}
\toprule
Operation & Folds noticed, at most\\
\midrule
Uniform entrywise function, transpose, permitted constant & $0$\\
Diagonal formation or extraction, broadcasting & $0$\\
Product over an index of size one & $0$\\
Vector inner product & $0$\\
Matrix--vector or vector--matrix contraction & $\beta$\\
Square matrix--matrix contraction & $3\beta$\\
\bottomrule
\end{tabular}
\end{center}
The cost counts noticed folds, not the time needed to evaluate a gate.

\paragraph{Why join breadth matters.}
The factor $\beta$ in the local matrix--vector bound cannot in
general be replaced by a constant. For $b\ge1$, the finite-set
semiring $(2^{[b]},\cup,\cap)$ has join breadth exactly $b$:
every union is determined by at most $b$ sets, whereas the union
of the $b$ singletons requires all of them. The same is true of $\mathsf{Trop}_{\mathbb Z}^{b}$, the $b$-fold
product of the integer min--plus semiring with coordinatewise
operations: different coordinates can attain their minima at
different summands.
Moreover, when $2^{k-1}\ge b$, choose distinct colors
$\lambda_1,\ldots,\lambda_b$ with no two complementary, and put
$X[\lambda_j]=\{j\}$ and $X[\mu]=\varnothing$ for all other
colors and on the diagonal. Matrix--vector aggregation of $X$
against the constant column $[b]$, using union for aggregation
and intersection as kernel, yields $[b]$ at the hub on the
reference structure. Folding $\lambda_j$ changes this value to
$[b]\setminus\{j\}$. Thus one aggregation notices $b$ folds.
This is local tightness, not a claim that the overall circuit bound
is sharp in $\beta$.

\subsection{Exponential arity-elimination lower bounds}\label{sec:arity-bounds}
We now apply Theorem~\ref{thm:test-cost} to the $2^k$ original fold
tests. The inputs are encoded by the two distinct semiring values
$0_S,1_S$. On every test, idempotence makes $\mathcal G_k$
return $1_S$ or $0_S$ according to the Boolean truth of $G_k$,
so each fold still changes the answer at the unchanged read edge.
For natural and integer min--plus the two values are $+\infty$
and $0$, respectively. Intermediate gates, however, may take
arbitrary values of $S$, which is why the scalar analysis of the
previous subsection is needed.

\tropicalmatrixbounds*
\begin{proof}
Both tropical semirings have join breadth $\beta=1$.
Use the $2^k$ fold tests of Lemma~\ref{lem:G-tests},
encoded as $0_S,1_S$. Theorem~\ref{thm:test-cost} and
Lemma~\ref{lem:typed-cost} give
$2^k\le\mvcount C+3\mmcount C$.
The ARA$(3)$ upper bound was constructed in \S\ref{sec:weightedG}.
\end{proof}

\begin{theorem}[Exponential arity-elimination overhead]\label{cor:matrix}
Let $S$ be a nontrivial commutative idempotent semiring of finite
join breadth $\beta$. For every $k\ge1$, the query $\mathcal G_k$
has a positive ARA$(3)$ expression of size $O(k)$, but every typed
matrix circuit $C$ computing it on all finite inputs satisfies
\[
 \beta(\mvcount C+3\mmcount C)\ge2^k.
\]
In particular, at least $2^k/(3\beta)$ aggregation gates are
necessary. Conversely, $2^k$ matrix--vector products suffice.
\end{theorem}
\begin{proof}
The $0_S,1_S$ encodings above supply $2^k$ fold tests.
Theorem~\ref{thm:test-cost} and Lemma~\ref{lem:typed-cost} give
the lower bound. Since
$\mvcount C+3\mmcount C\le3(\mvcount C+\mmcount C)$,
there must be at least $2^k/(3\beta)$ aggregation gates.

For the upper bound, the explicit MATLANG expression
\eqref{eq:matlang-dnf} uses exactly $2^k$ matrix--vector
aggregations; all remaining operations are entrywise or broadcasts.
Its correctness for arbitrary semiring-valued inputs follows from
distributivity, as verified in \S\ref{sec:weightedG}.
Reusing partial products in the $H_A$ as in
\S\ref{sec:main-results} gives $O(2^k)$ total nodes.
\end{proof}

\begin{theorem}[Boolean arity-elimination overhead]\label{thm:boolean-matrix}
Over the Boolean semiring $\mathbb B$, every typed matrix circuit
$C$ computing $G_k$ on all finite inputs satisfies
$\mvcount C+2\mmcount C\ge2^k$.
This holds with arbitrary uniform entrywise functions and
aggregation kernels.
\end{theorem}
\begin{proof}
The $2^k$ Boolean tests, Theorem~\ref{thm:test-cost}, and the
Boolean refinement of Lemma~\ref{lem:typed-cost} give the claim.
\end{proof}

\paragraph{Why the Boolean case is not just a restatement of the CoR bound.}
Boolean matrix multiplication is relational composition, so total-size
lower bounds for ordinary Boolean matrix expressions are closely
related to Theorem~\ref{thm:main}. Here we count only aggregations:
all other permitted matrix operations are uncharged in this measure. For example,
diagonal extraction followed by broadcasting obtains the matrix
$(u,v)\mapsto R(u,u)$ without aggregation, whereas a CoR expression
for it is $(R\cap\id);\one$, which uses a composition.
Consequently, the aggregation count of a matrix circuit cannot be
identified with the composition count of its CoR translation.
We instead bound the fold sensitivity of the typed operations directly;
that analysis also works over non-Boolean semirings.

Examples of breadth-one semirings include
$\mathsf{Trop}_{\mathbb N}$, $\mathsf{Trop}_{\mathbb Z}$,
the integer max--plus semiring
$(\mathbb Z\cup\{-\infty\},\max,+)$, and max--min on $[0,1]$.
All require at least $2^k/3$ aggregations, even in our enlarged
kernel model.

\subsection{Positive computations without idempotence}\label{sec:positive}
A separate transfer handles positive arithmetic, for which join breadth
is not available. Let $(S,+,\otimes,0_S,1_S)$ be a nontrivial commutative
semiring satisfying
\begin{equation*}\tag{NV}\label{eq:positive-characteristic}
 r\,1_S\ne0_S\qquad\text{for every integer }r\ge1,
\end{equation*}
where $r\,1_S$ is the sum of $r$ copies of $1_S$.
Restrict the target to positive entrywise addition and multiplication,
semiring matrix products, transpose, broadcasting, diagonal operations,
and constants in $\mathbb N1_S\coloneqq\{r\,1_S:r\ge0\}$.
Arbitrary entrywise decoding and negative constants are not allowed here.
Replace joins in $\mathcal G_k$ by sums to define its positive
semiring version.

On Boolean-encoded inputs, every intermediate belongs to $\mathbb N1_S$.
The map sending $0_S$ to zero and every other element of this subsemiring
to one is a homomorphism to the Boolean semiring. Indeed, sums and products
of positive multiples of $1_S$ are again nonzero by the nonvanishing
condition~\eqref{eq:positive-characteristic}. Applying this map to every gate of
a target circuit $C$ gives a Boolean matrix circuit with exactly the same
typed product counts.
Its output is $G_k$. The Boolean lower bound therefore gives
$\mvcount{C}+2\mmcount{C}\ge2^k$.
Here the cost is applied to the Booleanized computation; no argument by
predictions is claimed for the original scalar circuit.
The $2^k$-product expansion still works by distributivity. This includes
positive computations over $\mathbb N$ and nonnegative reals.

The restriction on the operations matters. Over $\mathbb N$, arbitrary
size-specific entrywise encoding and decoding can compute the Boolean
query $G_k$ on Boolean-encoded inputs using a single row-sum aggregation.
Indeed, assign each pair profile $p\in\{0,1\}^{2k}$ a distinct exponent
$j(p)\in\{0,\ldots,2^{2k}-1\}$. The sum
\[
 s_x\coloneqq\sum_{z\in V}(n+1)^{j(\pi(x,z))}
\]
encodes every profile count as a base-$(n+1)$ digit, with no carry
because each count is at most $n$. Broadcasting $s_x$ and decoding it
together with the read profile $\pi(x,y)$ determines $G_k(x,y)$.
Nonidempotence alone does not guarantee this encoding.
The finite-join-breadth assumption in Theorem~\ref{cor:matrix} prevents
one aggregation from depending on arbitrarily many indispensable colors.
